\clearpage
\section*{Additional references for the retained analytic note}
These two sources are already named in Section S17. Their bibliographic details
and consulted-version locators are supplied here so that the consolidated
package identifies them independently of the earlier A-line manuscript.
The note's other named sources---Ehm and colleagues, Giacomini and White, and
the New York Fed---appear in the main manuscript's References.

\paragraph{Learning-curve definitions.}
Viering, T., \& Loog, M. (2023). The shape of learning curves: A review.
\emph{IEEE Transactions on Pattern Analysis and Machine Intelligence},
\emph{45}(6), 7799--7819. \url{https://doi.org/10.1109/TPAMI.2022.3220744}.
The Section 2.1 locator in S17 refers to the consulted author version,
\href{https://arxiv.org/abs/2103.10948v2}{arXiv:2103.10948v2}, revised
5 November 2022, PDF page 2. Its journal issue year and author-version revision
date describe different version events. The cited discussion distinguishes a
fixed learning problem from averaging across problems; it does not validate
this study's history coordinate or a universal learning law.

\paragraph{Rearrangement over the probability index.}
Chernozhukov, V., Fern\'andez-Val, I., \& Galichon, A. (2010).
Quantile and probability curves without crossing.
\emph{Econometrica}, \emph{78}(3), 1093--1125.
\url{https://doi.org/10.3982/ECTA7880}.
The Proposition 4 locator in S17 refers to the consulted author version,
\href{https://arxiv.org/abs/0704.3649v3}{arXiv:0704.3649v3}, revised
15 July 2014, PDF page 12. Its \(L^p\) comparison concerns estimation error
of a quantile curve rearranged over its probability index. It is distinct from
the fixed-level, cross-horizon aggregate pinball identity proved in S17;
no identity between those results or new general rearrangement theorem is claimed.
